\uuid{xswJ}
\titre{Probabilités, événements}
\auteur{worms}
\organisation{exo7}
\licence{CC BY-NC-SA 4.0}
\datecreate{2017-08-28}
\niveau{L2}
\module{Probabilité et statistique}
\chapitre{Probabilité discrète}
\sousChapitre{Probabilité et dénombrement}
\difficulte{}
\theme{Probabilités, événements}
\exo7qcmid{worms.worms.019}

\contenu{
\begin{question}
\qclassification{260.01}
\qkeeporder
Soit $n$ un entier non nul et $k$ un entier compris entre $1$ et $n$. On considère un tableau contenant $n$ cases vides. Quel est le nombre de façons différentes de noircir $k$ de ces cases ?
\begin{answers}
\bad{$k$}
\bad{$n-k$}
\bad{$1-n$}
\good{Aucune des réponses précédentes.}
\end{answers}
\begin{explanations}
La réponse est évidemment $C_n^k$. Trop facile pour avoir faux !
\end{explanations}
\end{question}
}
